% GENERATED FILE. Edit canonical lessons, then run npm run generate:books.
% canonical-source-hash: fnv1a64:db67b926775d2e6d
% canonical-lessons: ZH-C382-zhiban, ZH-C382-baoan, ZH-W382-bao-keep, ZH-C382-jiankong, ZH-C382-baozheng, ZH-C382-xunluo

\chapter{The Security Guard}
\label{ch:zh-the-security-guard}
\hlchaptermodality{382}

\begin{chapteropening}
\textbf{By the end of this chapter:} \emph{I can talk about being on duty, security guards, cameras, promising and patrols.}

Complete the last lesson of chapter 382: I can talk about being on duty, security guards, cameras, promising and patrols.
\end{chapteropening}

\section[zhíbān]{zhíbān — to be on duty}

\label{lesson:ZH-C382-zhiban}



\begin{quote}
Before the new one: say the Mandarin for nationwide, then the Mandarin for a sense of security.
\end{quote}

\subsection*{You'll want to know: zhíbān}
\emph{zhíbān} — \textquotedblleft{}to be on duty\textquotedblright{}.

\begin{quote}
\emph{jīntiān wǎnshang wǒ zhíbān} — I'm on duty tonight.
\end{quote}

Its characters come later; for now it is a word for the ear.

\begin{cousinweb}
\emph{zhíbān} is built from \emph{zhí} 'to meet, to hold' and \emph{bān} 'a turn', both inherited from Middle Chinese, \emph{driH} and \emph{paen} (Guangyun, in Baxter's transcription), which glosses the first 'to hold' and 'to meet with' and the second 'to hand out tokens'. The sense 'a shift' is not given there.
\end{cousinweb}

\begin{grammarlens}[title={What you've built}]
One more word for giving reasons.
\end{grammarlens}

\subsection*{Your turn}
\begin{itemize}
\raggedright
  \item \emph{Say it:} \emph{zhíbān}
  \item \emph{Say it:} \emph{zhíbān}, once more
  \item \emph{Say it:} \emph{jīntiān wǎnshang wǒ zhíbān} — I'm on duty tonight.
  \item \emph{Recall:} say the Mandarin for nationwide, then the Mandarin for a sense of security, then say \emph{zhíbān} again
\end{itemize}

\begin{tcolorbox}[breakable,colback=teal!4,colframe=teal!35!black,title={Before you move on}]
What does zhíbān mean? (\textquotedblleft{}To be on duty\textquotedblright{}.) Say it once more.
\end{tcolorbox}
\section[\texorpdfstring{\zh{保安} (bǎo'ān)}{保安 (bǎo'ān)}]{\zh{保安} (bǎo'ān) — security guard}

\label{lesson:ZH-C382-baoan}



\begin{quote}
Before the new one: say the Mandarin for a sense of security, then the Mandarin for to be on duty.
\end{quote}

\subsection*{You'll want to know: \zh{保安}}
\textbf{\zh{保安}} — \emph{bǎo'ān} — \textquotedblleft{}security guard\textquotedblright{}.

\begin{quote}
\textbf{\zh{小区门口有两个保安}} \emph{xiǎoqū ménkǒu yǒu liǎng ge bǎo'ān} — There are two security guards at the gate of the estate.
\end{quote}

The new character, in pieces: \zh{保} — nine strokes in three pieces: the person piece \zh{亻}, then \zh{口} high on the right over a tree of four strokes. The next lesson puts it in your hand.

\begin{cousinweb}
\textbf{\zh{保安}} (\emph{bǎo'ān}) is built from \zh{保} 'to keep' and \zh{安} 'calm', both inherited from Middle Chinese, \emph{pawX} and \emph{'an} (Guangyun, in Baxter's transcription), which glosses the first 'to guard' and 'to keep safe': keeping the peace, and then the person who does it.
\end{cousinweb}

\begin{grammarlens}[title={What you've built}]
One more word for giving reasons.
\end{grammarlens}

\subsection*{Your turn}
\begin{itemize}
\raggedright
  \item \emph{Say it:} \emph{bǎo'ān}
  \item \emph{Say it:} \emph{bǎo'ān}, once more
  \item \emph{Say it:} \emph{xiǎoqū ménkǒu yǒu liǎng ge bǎo'ān} — There are two security guards at the gate of the estate.
  \item \emph{Recall:} say the Mandarin for a sense of security, then the Mandarin for to be on duty, then say \emph{bǎo'ān} again
\end{itemize}

\begin{tcolorbox}[breakable,colback=teal!4,colframe=teal!35!black,title={Before you move on}]
What does \zh{保安} mean? (\textquotedblleft{}Security guard\textquotedblright{}.) Say it once more.
\end{tcolorbox}
\section[\texorpdfstring{\zh{保} (bǎo)}{保 (bǎo)}]{\zh{保} — the person piece, then \zh{口} over a tree}

\label{lesson:ZH-W382-bao-keep}



\begin{quote}
The word first, then the character inside it.

\begin{itemize}
\raggedright
  \item \emph{Recall:} say \emph{bǎo'ān} — \textbf{R1}, one lesson back, before it decays
\end{itemize}

\begin{quote}
\zh{你} \zh{安}
\end{quote}

\emph{Read it:} each one aloud, with nothing said first — \emph{nǐ}, \emph{ān}
\end{quote}

\subsection*{Script — the person piece, then \zh{口} over a tree}
\hlblockfigure{\detokenize{figures/ZH-W382-bao-keep-filmstrip.pdf}}{How \zh{保} is written, stroke by stroke}

\begin{quote}
\zh{保}
\end{quote}

\textbf{Three pieces: the person piece on the left, then \zh{口} over a tree on the right.}

\textbf{The person piece \zh{亻}, two strokes,} as in \zh{你} and \zh{他}.

\textbf{\zh{口}, three strokes, high on the right:} a short vertical, a stroke that runs right and turns down, then the bar that closes the foot.

\textbf{The tree under it, four strokes.} A long bar. A vertical down through it. A stroke falling to the left from the crossing, then one falling to the right.

2 + 3 + 4 = nine.

\subsection*{Your turn}
\begin{itemize}
\raggedright
  \item \emph{Write it:} \zh{保}, all nine, with the model in front of you
  \item \emph{Cover:} write it again from memory, and count the pieces, not the strokes
\end{itemize}

\begin{tcolorbox}[breakable,colback=teal!4,colframe=teal!35!black,title={Before you move on}]
How many strokes in \zh{保}? (\textbf{Nine}.) What sits over the tree? (\textbf{\zh{口}}.) Which piece comes first? (\textbf{\zh{亻}, the person piece}.)

Source: \href{https://raw.githubusercontent.com/chanind/hanzi-writer-data/68d10a4b21150cae5e1ebbd223eed289cf32d90c/data/\%E4\%BF\%9D.json}{Hanzi Writer Data \zh{保}, pinned stroke paths and medians}.
\end{tcolorbox}
\section[jiānkòng]{jiānkòng — security camera; CCTV}

\label{lesson:ZH-C382-jiankong}



\begin{quote}
Before the new one: say the Mandarin for to be on duty, then the Mandarin for a security guard.
\end{quote}

\subsection*{You'll want to know: jiānkòng}
\emph{jiānkòng} — \textquotedblleft{}security camera; CCTV\textquotedblright{}.

\begin{quote}
\emph{yīnwèi yǒu jiānkòng suǒyǐ xiǎoqū lǐ hěn ānquán} — There are cameras, so the estate is safe.
\end{quote}

Its characters come later; for now it is a word for the ear.

\begin{cousinweb}
\emph{jiānkòng} is built from \emph{jiān} 'to watch over' and \emph{kòng} 'to pull', both inherited from Middle Chinese, \emph{kaem} and \emph{khuwngH} (Guangyun, in Baxter's transcription), which glosses the first 'to look down on and inspect' and the second 'to draw' and 'to report'. The camera is a modern sense.
\end{cousinweb}

\begin{grammarlens}[title={What you've built}]
One more word for giving reasons.
\end{grammarlens}

\subsection*{Your turn}
\begin{itemize}
\raggedright
  \item \emph{Say it:} \emph{jiānkòng}
  \item \emph{Say it:} \emph{jiānkòng}, once more
  \item \emph{Say it:} \emph{yīnwèi yǒu jiānkòng suǒyǐ xiǎoqū lǐ hěn ānquán} — There are cameras, so the estate is safe.
  \item \emph{Recall:} say the Mandarin for to be on duty, then the Mandarin for a security guard, then say \emph{jiānkòng} again
\end{itemize}

\begin{tcolorbox}[breakable,colback=teal!4,colframe=teal!35!black,title={Before you move on}]
What does jiānkòng mean? (\textquotedblleft{}Security camera; CCTV\textquotedblright{}.) Say it once more.
\end{tcolorbox}
\section[bǎozhèng]{bǎozhèng — to promise; to guarantee}

\label{lesson:ZH-C382-baozheng}



\begin{quote}
Before the new one: say the Mandarin for a security guard, then the Mandarin for a security camera.
\end{quote}

\subsection*{You'll want to know: bǎozhèng}
\emph{bǎozhèng} — \textquotedblleft{}to promise; to guarantee\textquotedblright{}.

\begin{quote}
\emph{wǒ bǎozhèng yǐhòu bú zhème chǎo le} — I promise I won't be this noisy again.
\end{quote}

Its characters come later; for now it is a word for the ear.

\begin{cousinweb}
\emph{bǎozhèng} is built from \emph{bǎo} 'to keep' and \emph{zhèng} 'to prove', both inherited from Middle Chinese, \emph{pawX} and \emph{tsyingH} (Guangyun, in Baxter's transcription), which glosses the first 'to answer for' and the second 'to test'.
\end{cousinweb}

\begin{grammarlens}[title={What you've built}]
One more word for giving reasons.
\end{grammarlens}

\subsection*{Your turn}
\begin{itemize}
\raggedright
  \item \emph{Say it:} \emph{bǎozhèng}
  \item \emph{Say it:} \emph{bǎozhèng}, once more
  \item \emph{Say it:} \emph{wǒ bǎozhèng yǐhòu bú zhème chǎo le} — I promise I won't be this noisy again.
  \item \emph{Recall:} say the Mandarin for a security guard, then the Mandarin for a security camera, then say \emph{bǎozhèng} again
\end{itemize}

\begin{tcolorbox}[breakable,colback=teal!4,colframe=teal!35!black,title={Before you move on}]
What does bǎozhèng mean? (\textquotedblleft{}To promise; to guarantee\textquotedblright{}.) Say it once more.
\end{tcolorbox}
\section[xúnluó]{xúnluó — to patrol}

\label{lesson:ZH-C382-xunluo}



\begin{quote}
Before the new one: say the Mandarin for a security camera, then the Mandarin for to guarantee.
\end{quote}

\subsection*{You'll want to know: xúnluó}
\emph{xúnluó} — \textquotedblleft{}to patrol\textquotedblright{}.

\begin{quote}
\emph{wǎnshang bǎo'ān yīnggāi zài xiǎoqū lǐ xúnluó} — At night the guards should patrol the estate.
\end{quote}

Its characters come later; for now it is a word for the ear.

\begin{cousinweb}
\emph{xúnluó} is built from \emph{xún} 'to go round' and \emph{luó} 'a patrol', both inherited from Middle Chinese, \emph{zwin} and \emph{laH} (Guangyun, in Baxter's transcription), which glosses the first 'to go about looking' and the second 'roving soldiers'. Mandarin luó has a rising tone where the dictionary gives a departing one.
\end{cousinweb}

\begin{grammarlens}[title={What you've built}]
One more word for giving reasons.
\end{grammarlens}

\subsection*{Your turn}
\begin{itemize}
\raggedright
  \item \emph{Say it:} \emph{xúnluó}
  \item \emph{Say it:} \emph{xúnluó}, once more
  \item \emph{Say it:} \emph{wǎnshang bǎo'ān yīnggāi zài xiǎoqū lǐ xúnluó} — At night the guards should patrol the estate.
  \item \emph{Recall:} say the Mandarin for a security camera, then the Mandarin for to guarantee, then say \emph{xúnluó} again
\end{itemize}

\begin{tcolorbox}[breakable,colback=teal!4,colframe=teal!35!black,title={Before you move on}]
What does xúnluó mean? (\textquotedblleft{}To patrol\textquotedblright{}.) Say it once more.
\end{tcolorbox}
